% Responsable: Martín (investigación)
% 10-15 fuentes. Cada subsección termina diciendo qué hueco llena este trabajo.
\section{Revisión de literatura}\label{sec:literatura}

\subsection{Medición multidimensional de la pobreza}
\pendiente{Alkire--Foster \citep{alkire2011}; IPM de Colombia \citep{angulo2016};
brecha normalizada / severidad.}

\subsection{Focalización y errores de inclusión y exclusión}
\pendiente{\citet{coady2004}; SISBÉN IV y focalización geográfica en Colombia.}

\subsection{Optimización aplicada a la reducción de pobreza}
\pendiente{\citet{hamie2024} y trabajos relacionados (MILP, asignación de presupuesto).}

\subsection{Tipologías de hogares mediante \emph{clustering}}
\pendiente{k-medoides / PAM \citep{kaufman1990}; uso en encuestas de hogares.}
\nota{Jerónimo}{Martín: esta subsección sostiene el Modelo 4. Si no alcanzas,
la escribo yo con lo de la elección de $k$.}
